% ============================================================
%  Chapter 6: Implementation
%  Included from report.tex via % ============================================================
%  Chapter 6: Implementation
%  Included from report.tex via % ============================================================
%  Chapter 6: Implementation
%  Included from report.tex via % ============================================================
%  Chapter 6: Implementation
%  Included from report.tex via \include{Chapters/Implementation}.
%  Figures fig-6-1 .. fig-6-14 are user screenshots; each is wrapped
%  in \IfFileExists so the document compiles before PNGs are added.
% ============================================================

\chapter{Implementation}
\label{chap:impl}

% Helper: show image if present, else visible placeholder box.
\newcommand{\implscreenshot}[2]{%
  \IfFileExists{#1}{%
    \includegraphics[width=\linewidth,keepaspectratio]{#1}}{%
    \fbox{\begin{minipage}{0.96\linewidth}\centering\footnotesize
    \texttt{#1}\\[0.35em]#2\end{minipage}}}}

% Pointer to future appendix material (B.x or C.x). Replace macro body
% when appendices are written so all chapter pointers update together.
\newcommand{\apxref}[1]{\textit{(see Appendix~#1 when added.)}}

% ------------------------------------------------------------
\section{From design to implementation}
\label{sec:impl-bridge}

Chapter~\ref{chap:design} fixed what must be built and why each
major choice beat its alternative (Figure~\ref{fig:design-decisions}).
This chapter shows where that design lives in code, what
non-trivial implementation choices were made while wiring it, and
where the running system diverged from the plan
(\S\ref{sec:impl-deviations}). Every row of
Figure~\ref{fig:design-decisions} maps to at least one module or
identifiable block listed in \S\ref{sec:impl-layout}; the four
hardest blocks (provider asymmetry, JSON hardening, null-aware
comparison, failure diagnosis) each get a screenshot in
\S\S\ref{sec:impl-provider}--\ref{sec:impl-evaldiag}.

% ------------------------------------------------------------
\section{Repository layout}
\label{sec:impl-layout}

The runtime is a single Python package under \texttt{src/} plus a
thin \texttt{cli.py} entry point. Together they are about
\texttt{2,700} lines. Evaluation, LLM adapters, and OCR sit in
subpackages so the pipeline imports a stable surface
(\texttt{DatasetLoader}, \texttt{get\_provider}, \texttt{run\_pipeline},
\texttt{evaluate\_run}, \texttt{compare\_runs},
\texttt{generate\_report}). The test tree under \texttt{tests/}
contains \texttt{531} test functions across twelve files
(\S\ref{sec:impl-testing}). Configuration and the zero-shot prompt
live outside the package at \texttt{config/default.yaml} and
\texttt{prompts/extraction\_v1.txt} so a run can be replayed from
artefacts without re-reading the source.

Figure~\ref{fig:impl-tree} evidences a single package root with no
parallel ``evaluation copy'' of the schema: the same
\texttt{BillExtraction} type is imported by the pipeline, both
providers, and the evaluator. Individual file purposes and line
counts are not readable from a tree alone.

\begin{figure}[ht]
    \centering
    \implscreenshot{Chapters/figures/fig-6-1-project-tree.png}{%
      Screenshot: project tree (\texttt{src/}, \texttt{tests/},
      \texttt{config/}, \texttt{prompts/}, \texttt{cli.py}).}
    \caption{Repository layout. Runtime code under \texttt{src/};
    tests mirror module names.}
    \label{fig:impl-tree}
\end{figure}

A per-module inventory (path, purpose, key symbols, approximate
line count) will be tabulated when the appendices are finalised.
\apxref{B.1}

% ------------------------------------------------------------
\section{Data contract}
\label{sec:impl-schema}

\texttt{BillExtraction} is a Pydantic v2 \texttt{BaseModel} with
twelve optional fields and a class-level \texttt{SCHEMA\_FIELDS}
list that fixes iteration order for evaluation and reporting.
\texttt{schema\_description()} builds the bullet list injected into
the prompt at \texttt{\{schema\_description\}}; the same strings
appear nowhere else, so a field rename is a single edit.

The first implementation attempt used Pydantic v1 because the
course template did not pin a version. OpenAI's Structured Outputs
call (\texttt{text\_format=BillExtraction}) failed to emit a valid
JSON schema from the model class. Pinning Pydantic
\texttt{\textgreater=2.0} in \texttt{requirements.txt} and
migrating the model to v2 metadata fixed the failure in one step.
The alternative, dropping Structured Outputs and parsing JSON by
hand for OpenAI too, was rejected because it would remove the
strongest guarantee on the modality comparison's OpenAI arm.

Figure~\ref{fig:impl-schema} shows only the canonical iteration order
\texttt{SCHEMA\_FIELDS}: the design point is that evaluation, reporting,
and the prompt builder all walk the same ordered list. The twelve
\texttt{Optional} field declarations and \texttt{schema\_description()}
live in the same file but are omitted here to keep the crop short.

\begin{figure}[ht]
    \centering
    \implscreenshot{Chapters/figures/fig-6-2-schema.png}{%
      Tight crop: \texttt{src/schema.py} lines 33--46
      (\texttt{SCHEMA\_FIELDS} list).}
    \caption{Canonical field order for iteration (\texttt{src/schema.py}).}
    \label{fig:impl-schema}
\end{figure}

\apxref{C.1}

% ------------------------------------------------------------
\section{Manifest-driven loading}
\label{sec:impl-loader}

\texttt{DatasetLoader.load\_and\_validate} chains six steps: read
CSV, assert required columns, reject duplicate
\texttt{document\_id}, validate cell values (language, utility type,
status, integer \texttt{page\_count}), filter to
\texttt{active} rows with \texttt{annotated} and \texttt{verified}
true, resolve PDF and ground-truth paths, then assert both files
exist. An early prototype collapsed validation into one function
that returned only the first error string. That made large
manifests painful to debug because a row with a bad language code
and a missing PDF looked like a single generic failure. Splitting
the chain keeps each failure mode attributable to one stage, which
is what reproducibility audits need when a row is rejected.

Figure~\ref{fig:impl-loader} shows the ordered call chain inside
\texttt{load\_and\_validate}: the design point is staged validation so
each failure mode maps to one step. Helper implementations
(\texttt{\_read\_csv}, \texttt{\_validate\_columns}, and the rest) are
not shown in the crop.

\begin{figure}[ht]
    \centering
    \implscreenshot{Chapters/figures/fig-6-3-loader.png}{%
      Tight crop: \texttt{src/dataset\_loader.py} lines 55--61
      (sequential validation chain).}
    \caption{Staged manifest validation (\texttt{src/dataset\_loader.py}).}
    \label{fig:impl-loader}
\end{figure}

\apxref{C.2}

% ------------------------------------------------------------
\section{Provider strategy in code}
\label{sec:impl-provider}

\subsection{Abstract contract}

\texttt{LLMProvider} defines \texttt{extract\_from\_text},
\texttt{extract\_from\_image}, and \texttt{get\_model\_id}. Both
extraction methods return the same dict shape so
\texttt{pipeline.py} never branches on provider name. The prompt is
always the rendered instruction; OCR text or image bytes are
passed as the user payload only, which keeps token counts stable
when the instruction is long.

Figure~\ref{fig:impl-llm-abc} shows the text-mode abstract signature and
its return-shape contract: the design point is that both modalities
must return the same dict keys so \texttt{pipeline.py} stays
provider-agnostic. The vision abstract and shared helpers are omitted
from the crop.

\begin{figure}[ht]
    \centering
    \implscreenshot{Chapters/figures/fig-6-4-llm-abc.png}{%
      Tight crop: \texttt{src/llm/base.py} lines 45--58
      (\texttt{extract\_from\_text} signature and Returns block).}
    \caption{Provider contract for text mode (\texttt{src/llm/base.py}).}
    \label{fig:impl-llm-abc}
\end{figure}

\apxref{C.3}

\subsection{OpenAI: Responses API and Structured Outputs}

\texttt{OpenAIProvider} calls \texttt{responses.parse} with
\texttt{text\_format=BillExtraction}. The API returns text that
already matches the schema; the pipeline still runs
\texttt{\_parse\_llm\_json} for uniformity, but the happy path is a
single \texttt{json.loads} on a clean string. This is the direct
code realisation of the ``Structured output handling'' row in
Figure~\ref{fig:design-decisions}.

Figure~\ref{fig:impl-openai-text} isolates the Structured Outputs call:
\texttt{instructions=prompt} and a user message carrying
\texttt{ocr\_text} split rules from data, matching
Chapter~\ref{chap:design}. Latency wrapping and result assembly are
not shown.

\begin{figure}[ht]
    \centering
    \implscreenshot{Chapters/figures/fig-6-5-openai-text.png}{%
      Tight crop: \texttt{src/llm/openai\_provider.py} lines 63--74
      (\texttt{responses.parse} with \texttt{text\_format=BillExtraction}).}
    \caption{OpenAI text mode with Structured Outputs
    (\texttt{src/llm/openai\_provider.py}).}
    \label{fig:impl-openai-text}
\end{figure}

\apxref{C.4}

\subsection{Anthropic: Messages API and fence stripping}

\texttt{AnthropicProvider} uses \texttt{messages.create} with the
same rendered prompt as the system message and OCR text or image
blocks as the user message. There is no Structured Outputs hook
for this schema, so \texttt{\_strip\_json\_fencing} removes
\texttt{```json ... ```} wrappers before the shared parser runs;
that helper is small enough to describe in prose only. Provider-specific
cleaning stops here; anything that survives a malformed fence still
hits the four-stage chain in \S\ref{sec:impl-pipeline}.

Figure~\ref{fig:impl-anthropic-text} shows the text-mode API call only:
\texttt{system=prompt} mirrors OpenAI's split of instruction versus
payload. Vision mode and \texttt{\_build\_result} are omitted from the
crop.

\begin{figure}[ht]
    \centering
    \implscreenshot{Chapters/figures/fig-6-6-anthropic-text.png}{%
      Tight crop: \texttt{src/llm/anthropic\_provider.py} lines 75--88
      (\texttt{messages.create} for text mode).}
    \caption{Anthropic text mode (\texttt{src/llm/anthropic\_provider.py}).}
    \label{fig:impl-anthropic-text}
\end{figure}

\apxref{C.5}

\subsection{Registry and lazy import}

\texttt{registry.py} maps string keys to dotted class paths and
imports the class only when \texttt{get\_provider} runs. That means
a host configured for OpenAI never imports \texttt{anthropic}, which
matters when only one SDK is installed. Adding a third provider is
one new module plus one line in \texttt{\_PROVIDERS}.

Figure~\ref{fig:impl-registry} shows only the string-to-class map: the
design point is that adding a provider is one dictionary entry plus
one module. \texttt{get\_provider}, lazy import, and constructor wiring
are left to the appendix listing.

\begin{figure}[ht]
    \centering
    \implscreenshot{Chapters/figures/fig-6-7-registry.png}{%
      Tight crop: \texttt{src/llm/registry.py} lines 16--19
      (\texttt{\_PROVIDERS}).}
    \caption{Provider registry map (\texttt{src/llm/registry.py}).}
    \label{fig:impl-registry}
\end{figure}

\apxref{C.6}

% ------------------------------------------------------------
\section{Pipeline orchestration}
\label{sec:impl-pipeline}

\texttt{\_process\_document} wraps each stage in
\texttt{Timer} from \texttt{src/performance.py}. The keys written to
\texttt{timings.json} (\texttt{pdf\_to\_images\_ms},
\texttt{ocr\_ms}, \texttt{llm\_call\_ms},
\texttt{parse\_normalise\_ms}, \texttt{total\_ms}) match the design
sequence diagram (Figure~\ref{fig:design-seq}). Vision mode uses a
zero-millisecond placeholder timer for OCR so downstream reporting
can still sum columns without branching on key presence.

Figure~\ref{fig:impl-pipeline-core} shows both arms of the modality
split: OCR-then-\texttt{extract\_from\_text} versus vision-then
\texttt{extract\_from\_image}, including the zero-millisecond OCR timer
placeholder in vision mode. Stage~4 parse, timings assembly, artefact
writes, and cost estimation sit below the crop.

\begin{figure}[ht]
    \centering
    \implscreenshot{Chapters/figures/fig-6-8-pipeline-core.png}{%
      Tight crop: \texttt{src/pipeline.py} lines 133--150
      (\texttt{if mode == "ocr\_text"} and \texttt{else} vision branch).}
    \caption{OCR and vision modality branches inside
    \texttt{\_process\_document} (\texttt{src/pipeline.py}).}
    \label{fig:impl-pipeline-core}
\end{figure}

\texttt{\_parse\_llm\_json} implements the four-stage hardening
chain from Figure~\ref{fig:design-parse}: direct
\texttt{json.loads}, fenced-block strip, brace-balanced substring
extraction, then \texttt{ValueError} with a 500-character excerpt.
The chain never calls the model again, which preserves the same
post-LLM processing path for every modality and provider, supporting
modality isolation (ER1).

Figure~\ref{fig:impl-json-parse} shows the docstring plus stage~1 only:
the design point is the ordered fallback chain (direct parse, fence
strip, brace-balanced substring, then a bounded \texttt{ValueError}
excerpt). Stages~2--4 are visible in the full function in the appendix.

\begin{figure}[ht]
    \centering
    \implscreenshot{Chapters/figures/fig-6-9-json-parse.png}{%
      Tight crop: \texttt{src/pipeline.py} lines 54--67
      (\texttt{\_parse\_llm\_json} docstring and first parse attempt).}
    \caption{JSON hardening contract and stage~1 (\texttt{src/pipeline.py}).}
    \label{fig:impl-json-parse}
\end{figure}

\texttt{run\_pipeline} snapshots \texttt{config.yaml}, the manifest,
and \texttt{prompt.txt} into the run directory, then iterates
documents with resume-by-default (\texttt{extraction.json} skip)
and per-document \texttt{try}/\texttt{except} that writes
\texttt{error.json} on failure. That is the mechanical realisation of
NFR1 (configuration and prompt snapshots), NFR3 (per-document error
isolation), and ER4 (artefacts sufficient to reconstruct the run)
from Chapter~\ref{chap:spec}.

\apxref{C.7}

% ------------------------------------------------------------
\section{Normalisation}
\label{sec:impl-normalisation}

\texttt{FIELD\_NORMALISERS} maps each schema field name to a
function. \texttt{normalise\_extraction} walks that map so adding a
field touches one Pydantic attribute and one table entry, not a
switch statement scattered across three files.

\texttt{\_parse\_numeric\_string} disambiguates European
(\texttt{1.234,56}) and US (\texttt{1,234.56}) formats by the
relative order of the last comma and dot, not by locale, because LLM
output does not carry locale metadata. That is a concrete
implementation bet: wrong guess on ambiguous \texttt{12,34} is
possible, but bills in the dataset rarely use two-digit
``thousands'' groups without a third group.

Figure~\ref{fig:impl-normalisers} shows the dispatch table only: the
design point is that new fields add one Pydantic attribute and one
table row instead of a scattered \texttt{if} chain. The EU/US numeric
parser and per-field normaliser bodies are described above and listed
in full in the appendix.

\begin{figure}[ht]
    \centering
    \implscreenshot{Chapters/figures/fig-6-10-normalisers.png}{%
      Tight crop: \texttt{src/normalisation.py} lines 317--332
      (\texttt{FIELD\_NORMALISERS}).}
    \caption{Dispatch-table normalisation (\texttt{src/normalisation.py}).}
    \label{fig:impl-normalisers}
\end{figure}

\apxref{C.8}

% ------------------------------------------------------------
\section{Evaluation and diagnosis}
\label{sec:impl-evaldiag}

\texttt{compare\_field} resolves the null 2$\times$2 table before
any value comparison, then applies type-specific rules: $\pm$0.01
tolerance on money and consumption floats, exact string match on
ISO dates after normalisation, exact match on other strings with a
Levenshtein ratio stored alongside for inspection only (pass or
fail remains exact match).

Figure~\ref{fig:impl-compare-field} shows \texttt{compare\_field} from
its signature through line~124: null flags, the full 2$\times$2 null
table, then the float and date arms and the opening of the string
\texttt{else}. The figure uses a compact layout for the null returns;
the string-comparison body, the \texttt{category} assignment, and the
final \texttt{return FieldResult} sit below the crop.

\begin{figure}[ht]
    \centering
    \implscreenshot{Chapters/figures/fig-6-11-compare-field.png}{%
      Tight crop: \texttt{src/evaluation/metrics.py} lines 70--124
      (\texttt{compare\_field}: null table and typed branches).}
    \caption{Null-aware field comparison (\texttt{src/evaluation/metrics.py}).}
    \label{fig:impl-compare-field}
\end{figure}

\apxref{C.9}

\texttt{diagnose\_document} reads \texttt{ocr\_text.txt} for
\texttt{ocr\_text} mode and applies the substring heuristic from
Chapter~\ref{chap:design}, which is the automated half of ER5
(failure attribution in OCR-mediated conditions). Vision mode labels
every non-correct field as \texttt{vision\_model\_failure}. The
implementation writes \texttt{diagnosis.json} only when the CLI passes
\texttt{--diagnose}, so batch runs stay cheap by default while FR6
remains available on demand.

Figure~\ref{fig:impl-diagnosis} shows the OCR-mode loop only: the
condition \texttt{fr.ground\_truth is None or \_value\_in\_text(...)}
means a missing ground-truth string skips substring search and is
labelled \texttt{llm\_extraction\_failure} (hallucination), while a
present ground truth is checked against \texttt{ocr\_text.txt} to
split OCR loss from LLM mis-read. Vision-mode diagnosis and JSON
serialisation are not in the crop.

\begin{figure}[ht]
    \centering
    \implscreenshot{Chapters/figures/fig-6-12-diagnosis.png}{%
      Tight crop: \texttt{src/evaluation/diagnosis.py} lines 62--73
      (OCR-mode \texttt{failure\_type} assignment).}
    \caption{OCR-mode failure attribution (\texttt{src/evaluation/diagnosis.py}).}
    \label{fig:impl-diagnosis}
\end{figure}

\apxref{C.10}

% ------------------------------------------------------------
\section{Deviations from the design}
\label{sec:impl-deviations}

Three deviations were anticipated in
\S\ref{sec:design-maintenance}; all three appear in the code paths
above. First, OpenAI and Anthropic use different APIs and different
output guarantees; that asymmetry is visible in
Figures~\ref{fig:impl-openai-text} and~\ref{fig:impl-anthropic-text}.
Second, wall-clock latency is always measured with \texttt{Timer}
around the pipeline stages; provider-reported latency fields are
not trusted for aggregation. Third, \texttt{TesseractOCR} catches
\texttt{TesseractError} and returns empty text so a missing
language pack degrades to an OCR-side failure mode that diagnosis
can label.

Two additional deviations surfaced during implementation. Fourth,
\texttt{compare\_field} records a Levenshtein ratio for string
fields but still classifies pass or fail on exact equality only.
The design text in Chapter~\ref{chap:design} mentioned similarity
for inspection; the code matches that: the ratio is persisted in
\texttt{evaluation.json} for analysis but does not change the
four-class label. Fifth, per-model USD pricing for
\texttt{estimate\_cost} lives in a hard-coded dict inside
\texttt{src/performance.py} rather than in \texttt{config.yaml}.
The trade is reproducibility from a single checkout: a marker who
checks out the same commit sees the same price table that produced
the run's \texttt{estimated\_cost\_usd}. Moving prices to YAML
would be a small refactor if billing rates need frequent edits.

None of these five items change the factorial structure, the null
taxonomy, or the modality-isolation boundary that the evaluation
chapter will test.

% ------------------------------------------------------------
\section{Testing}
\label{sec:impl-testing}

Tests are organised one file per module, with class names that read
as behaviour contracts (\texttt{TestNormaliseDateShouldParseKnownFormats},
\texttt{TestMissingColumnsShouldRaise}, and similar).
There are \texttt{531} \texttt{def test\_*} functions across the
twelve files; the densest files are
\texttt{test\_evaluation.py}, \texttt{test\_pipeline.py},
\texttt{test\_normalisation.py}, and \texttt{test\_dataset\_loader.py}.

Figure~\ref{fig:impl-tests} is a two-panel composite: the left crop
shows the parametrised date table (positive coverage in one place);
the right crop shows assertions on \texttt{error.json} after bad LLM
JSON (pipeline resilience). Other test modules are referenced only in
prose.

\begin{figure}[ht]
    \centering
    \implscreenshot{Chapters/figures/fig-6-13-tests.png}{%
      Two crops: \texttt{tests/test\_normalisation.py} lines 34--45;
      \texttt{tests/test\_pipeline.py} lines 576--582.}
    \caption{Representative tests: parametrised dates and
    \texttt{error.json} on parse failure.}
    \label{fig:impl-tests}
\end{figure}

\apxref{C.11}

% ------------------------------------------------------------
\section{Running a condition end to end}
\label{sec:impl-running}

The surface area is four Click commands:
\texttt{run}, \texttt{evaluate}, \texttt{compare}, \texttt{report}.
\texttt{run} deep-merges CLI overrides into YAML, executes
\texttt{run\_pipeline}, and prints the run directory path.
\texttt{evaluate} writes \texttt{evaluation.json};
\texttt{--diagnose} adds \texttt{diagnosis.json} per document.
\texttt{compare} builds \texttt{comparison.json} for the four cells
of the $2 \times 2$ design. \texttt{report} reads
\texttt{comparison.json} and writes the six chart PNGs plus
\texttt{summary.txt} described in Chapter~\ref{chap:spec}.

Figure~\ref{fig:impl-run-cli} shows terminal output from
\texttt{python cli.py run ...} or, if API keys are not configured,
\texttt{python cli.py --help}. The design point is a single Click group
with four subcommands; option wiring and imports are not visible in a
terminal crop.

\begin{figure}[ht]
    \centering
    \implscreenshot{Chapters/figures/fig-6-14-run-cli.png}{%
      Screenshot: terminal, \texttt{python cli.py run ...} (or
      \texttt{python cli.py --help}).}
    \caption{CLI execution: run pipeline or command help
    (\texttt{cli.py}).}
    \label{fig:impl-run-cli}
\end{figure}

\apxref{C.12}

% ------------------------------------------------------------
\section{Dependency and maintenance recap}
\label{sec:impl-maintenance}

Runtime Python dependencies are listed in \texttt{requirements.txt}:
\texttt{pydantic}, \texttt{pyyaml}, \texttt{pdf2image},
\texttt{Pillow}, \texttt{pytesseract}, \texttt{openai},
\texttt{anthropic}, \texttt{python-Levenshtein}, \texttt{click},
\texttt{matplotlib}, \texttt{seaborn}, \texttt{pandas},
\texttt{psutil}, \texttt{pytest}, \texttt{python-dotenv}. System
binaries still required on the host are Poppler (for
\texttt{pdf2image}) and Tesseract with language packs
\texttt{eng}, \texttt{deu}, \texttt{fra}, \texttt{ita}. Extensibility
remains the three one-touch additions from
Chapter~\ref{chap:design}: new provider module plus registry line;
new \texttt{OCREngine} subclass; new schema field plus normaliser
entry.

% ------------------------------------------------------------
\section{Summary and forward bridge}
\label{sec:impl-summary}

This chapter mapped the ten design decisions of
Figure~\ref{fig:design-decisions} to concrete modules and showed the
five bounded deviations in \S\ref{sec:impl-deviations}. The
implementation is about \texttt{2,700} lines of runtime Python and
\texttt{531} tests, with disk artefacts and CLI commands
that make a full condition reproducible from a checkout plus a run
directory snapshot.

Two experimental requirements from Chapter~\ref{chap:spec} anchor
the narrative above. ER1 (modality isolation: same reasoning engine,
prompt template, and post-processing across modalities within a
provider row) is realised by the single shared
\texttt{\_process\_document} path after rasterisation
(\S\ref{sec:impl-pipeline}), the shared \texttt{\_parse\_llm\_json}
and normalisation stages, and the identical \texttt{BillExtraction}
contract for both providers (\S\S\ref{sec:impl-schema}--\ref{sec:impl-normalisation}).
ER4 (reproducibility: every run archives configuration, manifest,
prompt, and per-document artefacts) is realised by
\texttt{run\_pipeline}'s directory snapshots and outputs
(\S\ref{sec:impl-pipeline}), which \texttt{tests/test\_pipeline.py}
exercises together with resume and error paths.

Chapter~\ref{chap:evaluation} executes the $2 \times 2$ experiment,
reports accuracy, cost, and latency by condition, slices results by
language and utility type, and returns to the rejected alternatives
in Figure~\ref{fig:design-decisions} to compare predicted outcomes
with observed ones.
.
%  Figures fig-6-1 .. fig-6-14 are user screenshots; each is wrapped
%  in \IfFileExists so the document compiles before PNGs are added.
% ============================================================

\chapter{Implementation}
\label{chap:impl}

% Helper: show image if present, else visible placeholder box.
\newcommand{\implscreenshot}[2]{%
  \IfFileExists{#1}{%
    \includegraphics[width=\linewidth,keepaspectratio]{#1}}{%
    \fbox{\begin{minipage}{0.96\linewidth}\centering\footnotesize
    \texttt{#1}\\[0.35em]#2\end{minipage}}}}

% Pointer to future appendix material (B.x or C.x). Replace macro body
% when appendices are written so all chapter pointers update together.
\newcommand{\apxref}[1]{\textit{(see Appendix~#1 when added.)}}

% ------------------------------------------------------------
\section{From design to implementation}
\label{sec:impl-bridge}

Chapter~\ref{chap:design} fixed what must be built and why each
major choice beat its alternative (Figure~\ref{fig:design-decisions}).
This chapter shows where that design lives in code, what
non-trivial implementation choices were made while wiring it, and
where the running system diverged from the plan
(\S\ref{sec:impl-deviations}). Every row of
Figure~\ref{fig:design-decisions} maps to at least one module or
identifiable block listed in \S\ref{sec:impl-layout}; the four
hardest blocks (provider asymmetry, JSON hardening, null-aware
comparison, failure diagnosis) each get a screenshot in
\S\S\ref{sec:impl-provider}--\ref{sec:impl-evaldiag}.

% ------------------------------------------------------------
\section{Repository layout}
\label{sec:impl-layout}

The runtime is a single Python package under \texttt{src/} plus a
thin \texttt{cli.py} entry point. Together they are about
\texttt{2,700} lines. Evaluation, LLM adapters, and OCR sit in
subpackages so the pipeline imports a stable surface
(\texttt{DatasetLoader}, \texttt{get\_provider}, \texttt{run\_pipeline},
\texttt{evaluate\_run}, \texttt{compare\_runs},
\texttt{generate\_report}). The test tree under \texttt{tests/}
contains \texttt{531} test functions across twelve files
(\S\ref{sec:impl-testing}). Configuration and the zero-shot prompt
live outside the package at \texttt{config/default.yaml} and
\texttt{prompts/extraction\_v1.txt} so a run can be replayed from
artefacts without re-reading the source.

Figure~\ref{fig:impl-tree} evidences a single package root with no
parallel ``evaluation copy'' of the schema: the same
\texttt{BillExtraction} type is imported by the pipeline, both
providers, and the evaluator. Individual file purposes and line
counts are not readable from a tree alone.

\begin{figure}[ht]
    \centering
    \implscreenshot{Chapters/figures/fig-6-1-project-tree.png}{%
      Screenshot: project tree (\texttt{src/}, \texttt{tests/},
      \texttt{config/}, \texttt{prompts/}, \texttt{cli.py}).}
    \caption{Repository layout. Runtime code under \texttt{src/};
    tests mirror module names.}
    \label{fig:impl-tree}
\end{figure}

A per-module inventory (path, purpose, key symbols, approximate
line count) will be tabulated when the appendices are finalised.
\apxref{B.1}

% ------------------------------------------------------------
\section{Data contract}
\label{sec:impl-schema}

\texttt{BillExtraction} is a Pydantic v2 \texttt{BaseModel} with
twelve optional fields and a class-level \texttt{SCHEMA\_FIELDS}
list that fixes iteration order for evaluation and reporting.
\texttt{schema\_description()} builds the bullet list injected into
the prompt at \texttt{\{schema\_description\}}; the same strings
appear nowhere else, so a field rename is a single edit.

The first implementation attempt used Pydantic v1 because the
course template did not pin a version. OpenAI's Structured Outputs
call (\texttt{text\_format=BillExtraction}) failed to emit a valid
JSON schema from the model class. Pinning Pydantic
\texttt{\textgreater=2.0} in \texttt{requirements.txt} and
migrating the model to v2 metadata fixed the failure in one step.
The alternative, dropping Structured Outputs and parsing JSON by
hand for OpenAI too, was rejected because it would remove the
strongest guarantee on the modality comparison's OpenAI arm.

Figure~\ref{fig:impl-schema} shows only the canonical iteration order
\texttt{SCHEMA\_FIELDS}: the design point is that evaluation, reporting,
and the prompt builder all walk the same ordered list. The twelve
\texttt{Optional} field declarations and \texttt{schema\_description()}
live in the same file but are omitted here to keep the crop short.

\begin{figure}[ht]
    \centering
    \implscreenshot{Chapters/figures/fig-6-2-schema.png}{%
      Tight crop: \texttt{src/schema.py} lines 33--46
      (\texttt{SCHEMA\_FIELDS} list).}
    \caption{Canonical field order for iteration (\texttt{src/schema.py}).}
    \label{fig:impl-schema}
\end{figure}

\apxref{C.1}

% ------------------------------------------------------------
\section{Manifest-driven loading}
\label{sec:impl-loader}

\texttt{DatasetLoader.load\_and\_validate} chains six steps: read
CSV, assert required columns, reject duplicate
\texttt{document\_id}, validate cell values (language, utility type,
status, integer \texttt{page\_count}), filter to
\texttt{active} rows with \texttt{annotated} and \texttt{verified}
true, resolve PDF and ground-truth paths, then assert both files
exist. An early prototype collapsed validation into one function
that returned only the first error string. That made large
manifests painful to debug because a row with a bad language code
and a missing PDF looked like a single generic failure. Splitting
the chain keeps each failure mode attributable to one stage, which
is what reproducibility audits need when a row is rejected.

Figure~\ref{fig:impl-loader} shows the ordered call chain inside
\texttt{load\_and\_validate}: the design point is staged validation so
each failure mode maps to one step. Helper implementations
(\texttt{\_read\_csv}, \texttt{\_validate\_columns}, and the rest) are
not shown in the crop.

\begin{figure}[ht]
    \centering
    \implscreenshot{Chapters/figures/fig-6-3-loader.png}{%
      Tight crop: \texttt{src/dataset\_loader.py} lines 55--61
      (sequential validation chain).}
    \caption{Staged manifest validation (\texttt{src/dataset\_loader.py}).}
    \label{fig:impl-loader}
\end{figure}

\apxref{C.2}

% ------------------------------------------------------------
\section{Provider strategy in code}
\label{sec:impl-provider}

\subsection{Abstract contract}

\texttt{LLMProvider} defines \texttt{extract\_from\_text},
\texttt{extract\_from\_image}, and \texttt{get\_model\_id}. Both
extraction methods return the same dict shape so
\texttt{pipeline.py} never branches on provider name. The prompt is
always the rendered instruction; OCR text or image bytes are
passed as the user payload only, which keeps token counts stable
when the instruction is long.

Figure~\ref{fig:impl-llm-abc} shows the text-mode abstract signature and
its return-shape contract: the design point is that both modalities
must return the same dict keys so \texttt{pipeline.py} stays
provider-agnostic. The vision abstract and shared helpers are omitted
from the crop.

\begin{figure}[ht]
    \centering
    \implscreenshot{Chapters/figures/fig-6-4-llm-abc.png}{%
      Tight crop: \texttt{src/llm/base.py} lines 45--58
      (\texttt{extract\_from\_text} signature and Returns block).}
    \caption{Provider contract for text mode (\texttt{src/llm/base.py}).}
    \label{fig:impl-llm-abc}
\end{figure}

\apxref{C.3}

\subsection{OpenAI: Responses API and Structured Outputs}

\texttt{OpenAIProvider} calls \texttt{responses.parse} with
\texttt{text\_format=BillExtraction}. The API returns text that
already matches the schema; the pipeline still runs
\texttt{\_parse\_llm\_json} for uniformity, but the happy path is a
single \texttt{json.loads} on a clean string. This is the direct
code realisation of the ``Structured output handling'' row in
Figure~\ref{fig:design-decisions}.

Figure~\ref{fig:impl-openai-text} isolates the Structured Outputs call:
\texttt{instructions=prompt} and a user message carrying
\texttt{ocr\_text} split rules from data, matching
Chapter~\ref{chap:design}. Latency wrapping and result assembly are
not shown.

\begin{figure}[ht]
    \centering
    \implscreenshot{Chapters/figures/fig-6-5-openai-text.png}{%
      Tight crop: \texttt{src/llm/openai\_provider.py} lines 63--74
      (\texttt{responses.parse} with \texttt{text\_format=BillExtraction}).}
    \caption{OpenAI text mode with Structured Outputs
    (\texttt{src/llm/openai\_provider.py}).}
    \label{fig:impl-openai-text}
\end{figure}

\apxref{C.4}

\subsection{Anthropic: Messages API and fence stripping}

\texttt{AnthropicProvider} uses \texttt{messages.create} with the
same rendered prompt as the system message and OCR text or image
blocks as the user message. There is no Structured Outputs hook
for this schema, so \texttt{\_strip\_json\_fencing} removes
\texttt{```json ... ```} wrappers before the shared parser runs;
that helper is small enough to describe in prose only. Provider-specific
cleaning stops here; anything that survives a malformed fence still
hits the four-stage chain in \S\ref{sec:impl-pipeline}.

Figure~\ref{fig:impl-anthropic-text} shows the text-mode API call only:
\texttt{system=prompt} mirrors OpenAI's split of instruction versus
payload. Vision mode and \texttt{\_build\_result} are omitted from the
crop.

\begin{figure}[ht]
    \centering
    \implscreenshot{Chapters/figures/fig-6-6-anthropic-text.png}{%
      Tight crop: \texttt{src/llm/anthropic\_provider.py} lines 75--88
      (\texttt{messages.create} for text mode).}
    \caption{Anthropic text mode (\texttt{src/llm/anthropic\_provider.py}).}
    \label{fig:impl-anthropic-text}
\end{figure}

\apxref{C.5}

\subsection{Registry and lazy import}

\texttt{registry.py} maps string keys to dotted class paths and
imports the class only when \texttt{get\_provider} runs. That means
a host configured for OpenAI never imports \texttt{anthropic}, which
matters when only one SDK is installed. Adding a third provider is
one new module plus one line in \texttt{\_PROVIDERS}.

Figure~\ref{fig:impl-registry} shows only the string-to-class map: the
design point is that adding a provider is one dictionary entry plus
one module. \texttt{get\_provider}, lazy import, and constructor wiring
are left to the appendix listing.

\begin{figure}[ht]
    \centering
    \implscreenshot{Chapters/figures/fig-6-7-registry.png}{%
      Tight crop: \texttt{src/llm/registry.py} lines 16--19
      (\texttt{\_PROVIDERS}).}
    \caption{Provider registry map (\texttt{src/llm/registry.py}).}
    \label{fig:impl-registry}
\end{figure}

\apxref{C.6}

% ------------------------------------------------------------
\section{Pipeline orchestration}
\label{sec:impl-pipeline}

\texttt{\_process\_document} wraps each stage in
\texttt{Timer} from \texttt{src/performance.py}. The keys written to
\texttt{timings.json} (\texttt{pdf\_to\_images\_ms},
\texttt{ocr\_ms}, \texttt{llm\_call\_ms},
\texttt{parse\_normalise\_ms}, \texttt{total\_ms}) match the design
sequence diagram (Figure~\ref{fig:design-seq}). Vision mode uses a
zero-millisecond placeholder timer for OCR so downstream reporting
can still sum columns without branching on key presence.

Figure~\ref{fig:impl-pipeline-core} shows both arms of the modality
split: OCR-then-\texttt{extract\_from\_text} versus vision-then
\texttt{extract\_from\_image}, including the zero-millisecond OCR timer
placeholder in vision mode. Stage~4 parse, timings assembly, artefact
writes, and cost estimation sit below the crop.

\begin{figure}[ht]
    \centering
    \implscreenshot{Chapters/figures/fig-6-8-pipeline-core.png}{%
      Tight crop: \texttt{src/pipeline.py} lines 133--150
      (\texttt{if mode == "ocr\_text"} and \texttt{else} vision branch).}
    \caption{OCR and vision modality branches inside
    \texttt{\_process\_document} (\texttt{src/pipeline.py}).}
    \label{fig:impl-pipeline-core}
\end{figure}

\texttt{\_parse\_llm\_json} implements the four-stage hardening
chain from Figure~\ref{fig:design-parse}: direct
\texttt{json.loads}, fenced-block strip, brace-balanced substring
extraction, then \texttt{ValueError} with a 500-character excerpt.
The chain never calls the model again, which preserves the same
post-LLM processing path for every modality and provider, supporting
modality isolation (ER1).

Figure~\ref{fig:impl-json-parse} shows the docstring plus stage~1 only:
the design point is the ordered fallback chain (direct parse, fence
strip, brace-balanced substring, then a bounded \texttt{ValueError}
excerpt). Stages~2--4 are visible in the full function in the appendix.

\begin{figure}[ht]
    \centering
    \implscreenshot{Chapters/figures/fig-6-9-json-parse.png}{%
      Tight crop: \texttt{src/pipeline.py} lines 54--67
      (\texttt{\_parse\_llm\_json} docstring and first parse attempt).}
    \caption{JSON hardening contract and stage~1 (\texttt{src/pipeline.py}).}
    \label{fig:impl-json-parse}
\end{figure}

\texttt{run\_pipeline} snapshots \texttt{config.yaml}, the manifest,
and \texttt{prompt.txt} into the run directory, then iterates
documents with resume-by-default (\texttt{extraction.json} skip)
and per-document \texttt{try}/\texttt{except} that writes
\texttt{error.json} on failure. That is the mechanical realisation of
NFR1 (configuration and prompt snapshots), NFR3 (per-document error
isolation), and ER4 (artefacts sufficient to reconstruct the run)
from Chapter~\ref{chap:spec}.

\apxref{C.7}

% ------------------------------------------------------------
\section{Normalisation}
\label{sec:impl-normalisation}

\texttt{FIELD\_NORMALISERS} maps each schema field name to a
function. \texttt{normalise\_extraction} walks that map so adding a
field touches one Pydantic attribute and one table entry, not a
switch statement scattered across three files.

\texttt{\_parse\_numeric\_string} disambiguates European
(\texttt{1.234,56}) and US (\texttt{1,234.56}) formats by the
relative order of the last comma and dot, not by locale, because LLM
output does not carry locale metadata. That is a concrete
implementation bet: wrong guess on ambiguous \texttt{12,34} is
possible, but bills in the dataset rarely use two-digit
``thousands'' groups without a third group.

Figure~\ref{fig:impl-normalisers} shows the dispatch table only: the
design point is that new fields add one Pydantic attribute and one
table row instead of a scattered \texttt{if} chain. The EU/US numeric
parser and per-field normaliser bodies are described above and listed
in full in the appendix.

\begin{figure}[ht]
    \centering
    \implscreenshot{Chapters/figures/fig-6-10-normalisers.png}{%
      Tight crop: \texttt{src/normalisation.py} lines 317--332
      (\texttt{FIELD\_NORMALISERS}).}
    \caption{Dispatch-table normalisation (\texttt{src/normalisation.py}).}
    \label{fig:impl-normalisers}
\end{figure}

\apxref{C.8}

% ------------------------------------------------------------
\section{Evaluation and diagnosis}
\label{sec:impl-evaldiag}

\texttt{compare\_field} resolves the null 2$\times$2 table before
any value comparison, then applies type-specific rules: $\pm$0.01
tolerance on money and consumption floats, exact string match on
ISO dates after normalisation, exact match on other strings with a
Levenshtein ratio stored alongside for inspection only (pass or
fail remains exact match).

Figure~\ref{fig:impl-compare-field} shows \texttt{compare\_field} from
its signature through line~124: null flags, the full 2$\times$2 null
table, then the float and date arms and the opening of the string
\texttt{else}. The figure uses a compact layout for the null returns;
the string-comparison body, the \texttt{category} assignment, and the
final \texttt{return FieldResult} sit below the crop.

\begin{figure}[ht]
    \centering
    \implscreenshot{Chapters/figures/fig-6-11-compare-field.png}{%
      Tight crop: \texttt{src/evaluation/metrics.py} lines 70--124
      (\texttt{compare\_field}: null table and typed branches).}
    \caption{Null-aware field comparison (\texttt{src/evaluation/metrics.py}).}
    \label{fig:impl-compare-field}
\end{figure}

\apxref{C.9}

\texttt{diagnose\_document} reads \texttt{ocr\_text.txt} for
\texttt{ocr\_text} mode and applies the substring heuristic from
Chapter~\ref{chap:design}, which is the automated half of ER5
(failure attribution in OCR-mediated conditions). Vision mode labels
every non-correct field as \texttt{vision\_model\_failure}. The
implementation writes \texttt{diagnosis.json} only when the CLI passes
\texttt{--diagnose}, so batch runs stay cheap by default while FR6
remains available on demand.

Figure~\ref{fig:impl-diagnosis} shows the OCR-mode loop only: the
condition \texttt{fr.ground\_truth is None or \_value\_in\_text(...)}
means a missing ground-truth string skips substring search and is
labelled \texttt{llm\_extraction\_failure} (hallucination), while a
present ground truth is checked against \texttt{ocr\_text.txt} to
split OCR loss from LLM mis-read. Vision-mode diagnosis and JSON
serialisation are not in the crop.

\begin{figure}[ht]
    \centering
    \implscreenshot{Chapters/figures/fig-6-12-diagnosis.png}{%
      Tight crop: \texttt{src/evaluation/diagnosis.py} lines 62--73
      (OCR-mode \texttt{failure\_type} assignment).}
    \caption{OCR-mode failure attribution (\texttt{src/evaluation/diagnosis.py}).}
    \label{fig:impl-diagnosis}
\end{figure}

\apxref{C.10}

% ------------------------------------------------------------
\section{Deviations from the design}
\label{sec:impl-deviations}

Three deviations were anticipated in
\S\ref{sec:design-maintenance}; all three appear in the code paths
above. First, OpenAI and Anthropic use different APIs and different
output guarantees; that asymmetry is visible in
Figures~\ref{fig:impl-openai-text} and~\ref{fig:impl-anthropic-text}.
Second, wall-clock latency is always measured with \texttt{Timer}
around the pipeline stages; provider-reported latency fields are
not trusted for aggregation. Third, \texttt{TesseractOCR} catches
\texttt{TesseractError} and returns empty text so a missing
language pack degrades to an OCR-side failure mode that diagnosis
can label.

Two additional deviations surfaced during implementation. Fourth,
\texttt{compare\_field} records a Levenshtein ratio for string
fields but still classifies pass or fail on exact equality only.
The design text in Chapter~\ref{chap:design} mentioned similarity
for inspection; the code matches that: the ratio is persisted in
\texttt{evaluation.json} for analysis but does not change the
four-class label. Fifth, per-model USD pricing for
\texttt{estimate\_cost} lives in a hard-coded dict inside
\texttt{src/performance.py} rather than in \texttt{config.yaml}.
The trade is reproducibility from a single checkout: a marker who
checks out the same commit sees the same price table that produced
the run's \texttt{estimated\_cost\_usd}. Moving prices to YAML
would be a small refactor if billing rates need frequent edits.

None of these five items change the factorial structure, the null
taxonomy, or the modality-isolation boundary that the evaluation
chapter will test.

% ------------------------------------------------------------
\section{Testing}
\label{sec:impl-testing}

Tests are organised one file per module, with class names that read
as behaviour contracts (\texttt{TestNormaliseDateShouldParseKnownFormats},
\texttt{TestMissingColumnsShouldRaise}, and similar).
There are \texttt{531} \texttt{def test\_*} functions across the
twelve files; the densest files are
\texttt{test\_evaluation.py}, \texttt{test\_pipeline.py},
\texttt{test\_normalisation.py}, and \texttt{test\_dataset\_loader.py}.

Figure~\ref{fig:impl-tests} is a two-panel composite: the left crop
shows the parametrised date table (positive coverage in one place);
the right crop shows assertions on \texttt{error.json} after bad LLM
JSON (pipeline resilience). Other test modules are referenced only in
prose.

\begin{figure}[ht]
    \centering
    \implscreenshot{Chapters/figures/fig-6-13-tests.png}{%
      Two crops: \texttt{tests/test\_normalisation.py} lines 34--45;
      \texttt{tests/test\_pipeline.py} lines 576--582.}
    \caption{Representative tests: parametrised dates and
    \texttt{error.json} on parse failure.}
    \label{fig:impl-tests}
\end{figure}

\apxref{C.11}

% ------------------------------------------------------------
\section{Running a condition end to end}
\label{sec:impl-running}

The surface area is four Click commands:
\texttt{run}, \texttt{evaluate}, \texttt{compare}, \texttt{report}.
\texttt{run} deep-merges CLI overrides into YAML, executes
\texttt{run\_pipeline}, and prints the run directory path.
\texttt{evaluate} writes \texttt{evaluation.json};
\texttt{--diagnose} adds \texttt{diagnosis.json} per document.
\texttt{compare} builds \texttt{comparison.json} for the four cells
of the $2 \times 2$ design. \texttt{report} reads
\texttt{comparison.json} and writes the six chart PNGs plus
\texttt{summary.txt} described in Chapter~\ref{chap:spec}.

Figure~\ref{fig:impl-run-cli} shows terminal output from
\texttt{python cli.py run ...} or, if API keys are not configured,
\texttt{python cli.py --help}. The design point is a single Click group
with four subcommands; option wiring and imports are not visible in a
terminal crop.

\begin{figure}[ht]
    \centering
    \implscreenshot{Chapters/figures/fig-6-14-run-cli.png}{%
      Screenshot: terminal, \texttt{python cli.py run ...} (or
      \texttt{python cli.py --help}).}
    \caption{CLI execution: run pipeline or command help
    (\texttt{cli.py}).}
    \label{fig:impl-run-cli}
\end{figure}

\apxref{C.12}

% ------------------------------------------------------------
\section{Dependency and maintenance recap}
\label{sec:impl-maintenance}

Runtime Python dependencies are listed in \texttt{requirements.txt}:
\texttt{pydantic}, \texttt{pyyaml}, \texttt{pdf2image},
\texttt{Pillow}, \texttt{pytesseract}, \texttt{openai},
\texttt{anthropic}, \texttt{python-Levenshtein}, \texttt{click},
\texttt{matplotlib}, \texttt{seaborn}, \texttt{pandas},
\texttt{psutil}, \texttt{pytest}, \texttt{python-dotenv}. System
binaries still required on the host are Poppler (for
\texttt{pdf2image}) and Tesseract with language packs
\texttt{eng}, \texttt{deu}, \texttt{fra}, \texttt{ita}. Extensibility
remains the three one-touch additions from
Chapter~\ref{chap:design}: new provider module plus registry line;
new \texttt{OCREngine} subclass; new schema field plus normaliser
entry.

% ------------------------------------------------------------
\section{Summary and forward bridge}
\label{sec:impl-summary}

This chapter mapped the ten design decisions of
Figure~\ref{fig:design-decisions} to concrete modules and showed the
five bounded deviations in \S\ref{sec:impl-deviations}. The
implementation is about \texttt{2,700} lines of runtime Python and
\texttt{531} tests, with disk artefacts and CLI commands
that make a full condition reproducible from a checkout plus a run
directory snapshot.

Two experimental requirements from Chapter~\ref{chap:spec} anchor
the narrative above. ER1 (modality isolation: same reasoning engine,
prompt template, and post-processing across modalities within a
provider row) is realised by the single shared
\texttt{\_process\_document} path after rasterisation
(\S\ref{sec:impl-pipeline}), the shared \texttt{\_parse\_llm\_json}
and normalisation stages, and the identical \texttt{BillExtraction}
contract for both providers (\S\S\ref{sec:impl-schema}--\ref{sec:impl-normalisation}).
ER4 (reproducibility: every run archives configuration, manifest,
prompt, and per-document artefacts) is realised by
\texttt{run\_pipeline}'s directory snapshots and outputs
(\S\ref{sec:impl-pipeline}), which \texttt{tests/test\_pipeline.py}
exercises together with resume and error paths.

Chapter~\ref{chap:evaluation} executes the $2 \times 2$ experiment,
reports accuracy, cost, and latency by condition, slices results by
language and utility type, and returns to the rejected alternatives
in Figure~\ref{fig:design-decisions} to compare predicted outcomes
with observed ones.
.
%  Figures fig-6-1 .. fig-6-14 are user screenshots; each is wrapped
%  in \IfFileExists so the document compiles before PNGs are added.
% ============================================================

\chapter{Implementation}
\label{chap:impl}

% Helper: show image if present, else visible placeholder box.
\newcommand{\implscreenshot}[2]{%
  \IfFileExists{#1}{%
    \includegraphics[width=\linewidth,keepaspectratio]{#1}}{%
    \fbox{\begin{minipage}{0.96\linewidth}\centering\footnotesize
    \texttt{#1}\\[0.35em]#2\end{minipage}}}}

% Pointer to future appendix material (B.x or C.x). Replace macro body
% when appendices are written so all chapter pointers update together.
\newcommand{\apxref}[1]{\textit{(see Appendix~#1 when added.)}}

% ------------------------------------------------------------
\section{From design to implementation}
\label{sec:impl-bridge}

Chapter~\ref{chap:design} fixed what must be built and why each
major choice beat its alternative (Figure~\ref{fig:design-decisions}).
This chapter shows where that design lives in code, what
non-trivial implementation choices were made while wiring it, and
where the running system diverged from the plan
(\S\ref{sec:impl-deviations}). Every row of
Figure~\ref{fig:design-decisions} maps to at least one module or
identifiable block listed in \S\ref{sec:impl-layout}; the four
hardest blocks (provider asymmetry, JSON hardening, null-aware
comparison, failure diagnosis) each get a screenshot in
\S\S\ref{sec:impl-provider}--\ref{sec:impl-evaldiag}.

% ------------------------------------------------------------
\section{Repository layout}
\label{sec:impl-layout}

The runtime is a single Python package under \texttt{src/} plus a
thin \texttt{cli.py} entry point. Together they are about
\texttt{2,700} lines. Evaluation, LLM adapters, and OCR sit in
subpackages so the pipeline imports a stable surface
(\texttt{DatasetLoader}, \texttt{get\_provider}, \texttt{run\_pipeline},
\texttt{evaluate\_run}, \texttt{compare\_runs},
\texttt{generate\_report}). The test tree under \texttt{tests/}
contains \texttt{531} test functions across twelve files
(\S\ref{sec:impl-testing}). Configuration and the zero-shot prompt
live outside the package at \texttt{config/default.yaml} and
\texttt{prompts/extraction\_v1.txt} so a run can be replayed from
artefacts without re-reading the source.

Figure~\ref{fig:impl-tree} evidences a single package root with no
parallel ``evaluation copy'' of the schema: the same
\texttt{BillExtraction} type is imported by the pipeline, both
providers, and the evaluator. Individual file purposes and line
counts are not readable from a tree alone.

\begin{figure}[ht]
    \centering
    \implscreenshot{Chapters/figures/fig-6-1-project-tree.png}{%
      Screenshot: project tree (\texttt{src/}, \texttt{tests/},
      \texttt{config/}, \texttt{prompts/}, \texttt{cli.py}).}
    \caption{Repository layout. Runtime code under \texttt{src/};
    tests mirror module names.}
    \label{fig:impl-tree}
\end{figure}

A per-module inventory (path, purpose, key symbols, approximate
line count) will be tabulated when the appendices are finalised.
\apxref{B.1}

% ------------------------------------------------------------
\section{Data contract}
\label{sec:impl-schema}

\texttt{BillExtraction} is a Pydantic v2 \texttt{BaseModel} with
twelve optional fields and a class-level \texttt{SCHEMA\_FIELDS}
list that fixes iteration order for evaluation and reporting.
\texttt{schema\_description()} builds the bullet list injected into
the prompt at \texttt{\{schema\_description\}}; the same strings
appear nowhere else, so a field rename is a single edit.

The first implementation attempt used Pydantic v1 because the
course template did not pin a version. OpenAI's Structured Outputs
call (\texttt{text\_format=BillExtraction}) failed to emit a valid
JSON schema from the model class. Pinning Pydantic
\texttt{\textgreater=2.0} in \texttt{requirements.txt} and
migrating the model to v2 metadata fixed the failure in one step.
The alternative, dropping Structured Outputs and parsing JSON by
hand for OpenAI too, was rejected because it would remove the
strongest guarantee on the modality comparison's OpenAI arm.

Figure~\ref{fig:impl-schema} shows only the canonical iteration order
\texttt{SCHEMA\_FIELDS}: the design point is that evaluation, reporting,
and the prompt builder all walk the same ordered list. The twelve
\texttt{Optional} field declarations and \texttt{schema\_description()}
live in the same file but are omitted here to keep the crop short.

\begin{figure}[ht]
    \centering
    \implscreenshot{Chapters/figures/fig-6-2-schema.png}{%
      Tight crop: \texttt{src/schema.py} lines 33--46
      (\texttt{SCHEMA\_FIELDS} list).}
    \caption{Canonical field order for iteration (\texttt{src/schema.py}).}
    \label{fig:impl-schema}
\end{figure}

\apxref{C.1}

% ------------------------------------------------------------
\section{Manifest-driven loading}
\label{sec:impl-loader}

\texttt{DatasetLoader.load\_and\_validate} chains six steps: read
CSV, assert required columns, reject duplicate
\texttt{document\_id}, validate cell values (language, utility type,
status, integer \texttt{page\_count}), filter to
\texttt{active} rows with \texttt{annotated} and \texttt{verified}
true, resolve PDF and ground-truth paths, then assert both files
exist. An early prototype collapsed validation into one function
that returned only the first error string. That made large
manifests painful to debug because a row with a bad language code
and a missing PDF looked like a single generic failure. Splitting
the chain keeps each failure mode attributable to one stage, which
is what reproducibility audits need when a row is rejected.

Figure~\ref{fig:impl-loader} shows the ordered call chain inside
\texttt{load\_and\_validate}: the design point is staged validation so
each failure mode maps to one step. Helper implementations
(\texttt{\_read\_csv}, \texttt{\_validate\_columns}, and the rest) are
not shown in the crop.

\begin{figure}[ht]
    \centering
    \implscreenshot{Chapters/figures/fig-6-3-loader.png}{%
      Tight crop: \texttt{src/dataset\_loader.py} lines 55--61
      (sequential validation chain).}
    \caption{Staged manifest validation (\texttt{src/dataset\_loader.py}).}
    \label{fig:impl-loader}
\end{figure}

\apxref{C.2}

% ------------------------------------------------------------
\section{Provider strategy in code}
\label{sec:impl-provider}

\subsection{Abstract contract}

\texttt{LLMProvider} defines \texttt{extract\_from\_text},
\texttt{extract\_from\_image}, and \texttt{get\_model\_id}. Both
extraction methods return the same dict shape so
\texttt{pipeline.py} never branches on provider name. The prompt is
always the rendered instruction; OCR text or image bytes are
passed as the user payload only, which keeps token counts stable
when the instruction is long.

Figure~\ref{fig:impl-llm-abc} shows the text-mode abstract signature and
its return-shape contract: the design point is that both modalities
must return the same dict keys so \texttt{pipeline.py} stays
provider-agnostic. The vision abstract and shared helpers are omitted
from the crop.

\begin{figure}[ht]
    \centering
    \implscreenshot{Chapters/figures/fig-6-4-llm-abc.png}{%
      Tight crop: \texttt{src/llm/base.py} lines 45--58
      (\texttt{extract\_from\_text} signature and Returns block).}
    \caption{Provider contract for text mode (\texttt{src/llm/base.py}).}
    \label{fig:impl-llm-abc}
\end{figure}

\apxref{C.3}

\subsection{OpenAI: Responses API and Structured Outputs}

\texttt{OpenAIProvider} calls \texttt{responses.parse} with
\texttt{text\_format=BillExtraction}. The API returns text that
already matches the schema; the pipeline still runs
\texttt{\_parse\_llm\_json} for uniformity, but the happy path is a
single \texttt{json.loads} on a clean string. This is the direct
code realisation of the ``Structured output handling'' row in
Figure~\ref{fig:design-decisions}.

Figure~\ref{fig:impl-openai-text} isolates the Structured Outputs call:
\texttt{instructions=prompt} and a user message carrying
\texttt{ocr\_text} split rules from data, matching
Chapter~\ref{chap:design}. Latency wrapping and result assembly are
not shown.

\begin{figure}[ht]
    \centering
    \implscreenshot{Chapters/figures/fig-6-5-openai-text.png}{%
      Tight crop: \texttt{src/llm/openai\_provider.py} lines 63--74
      (\texttt{responses.parse} with \texttt{text\_format=BillExtraction}).}
    \caption{OpenAI text mode with Structured Outputs
    (\texttt{src/llm/openai\_provider.py}).}
    \label{fig:impl-openai-text}
\end{figure}

\apxref{C.4}

\subsection{Anthropic: Messages API and fence stripping}

\texttt{AnthropicProvider} uses \texttt{messages.create} with the
same rendered prompt as the system message and OCR text or image
blocks as the user message. There is no Structured Outputs hook
for this schema, so \texttt{\_strip\_json\_fencing} removes
\texttt{```json ... ```} wrappers before the shared parser runs;
that helper is small enough to describe in prose only. Provider-specific
cleaning stops here; anything that survives a malformed fence still
hits the four-stage chain in \S\ref{sec:impl-pipeline}.

Figure~\ref{fig:impl-anthropic-text} shows the text-mode API call only:
\texttt{system=prompt} mirrors OpenAI's split of instruction versus
payload. Vision mode and \texttt{\_build\_result} are omitted from the
crop.

\begin{figure}[ht]
    \centering
    \implscreenshot{Chapters/figures/fig-6-6-anthropic-text.png}{%
      Tight crop: \texttt{src/llm/anthropic\_provider.py} lines 75--88
      (\texttt{messages.create} for text mode).}
    \caption{Anthropic text mode (\texttt{src/llm/anthropic\_provider.py}).}
    \label{fig:impl-anthropic-text}
\end{figure}

\apxref{C.5}

\subsection{Registry and lazy import}

\texttt{registry.py} maps string keys to dotted class paths and
imports the class only when \texttt{get\_provider} runs. That means
a host configured for OpenAI never imports \texttt{anthropic}, which
matters when only one SDK is installed. Adding a third provider is
one new module plus one line in \texttt{\_PROVIDERS}.

Figure~\ref{fig:impl-registry} shows only the string-to-class map: the
design point is that adding a provider is one dictionary entry plus
one module. \texttt{get\_provider}, lazy import, and constructor wiring
are left to the appendix listing.

\begin{figure}[ht]
    \centering
    \implscreenshot{Chapters/figures/fig-6-7-registry.png}{%
      Tight crop: \texttt{src/llm/registry.py} lines 16--19
      (\texttt{\_PROVIDERS}).}
    \caption{Provider registry map (\texttt{src/llm/registry.py}).}
    \label{fig:impl-registry}
\end{figure}

\apxref{C.6}

% ------------------------------------------------------------
\section{Pipeline orchestration}
\label{sec:impl-pipeline}

\texttt{\_process\_document} wraps each stage in
\texttt{Timer} from \texttt{src/performance.py}. The keys written to
\texttt{timings.json} (\texttt{pdf\_to\_images\_ms},
\texttt{ocr\_ms}, \texttt{llm\_call\_ms},
\texttt{parse\_normalise\_ms}, \texttt{total\_ms}) match the design
sequence diagram (Figure~\ref{fig:design-seq}). Vision mode uses a
zero-millisecond placeholder timer for OCR so downstream reporting
can still sum columns without branching on key presence.

Figure~\ref{fig:impl-pipeline-core} shows both arms of the modality
split: OCR-then-\texttt{extract\_from\_text} versus vision-then
\texttt{extract\_from\_image}, including the zero-millisecond OCR timer
placeholder in vision mode. Stage~4 parse, timings assembly, artefact
writes, and cost estimation sit below the crop.

\begin{figure}[ht]
    \centering
    \implscreenshot{Chapters/figures/fig-6-8-pipeline-core.png}{%
      Tight crop: \texttt{src/pipeline.py} lines 133--150
      (\texttt{if mode == "ocr\_text"} and \texttt{else} vision branch).}
    \caption{OCR and vision modality branches inside
    \texttt{\_process\_document} (\texttt{src/pipeline.py}).}
    \label{fig:impl-pipeline-core}
\end{figure}

\texttt{\_parse\_llm\_json} implements the four-stage hardening
chain from Figure~\ref{fig:design-parse}: direct
\texttt{json.loads}, fenced-block strip, brace-balanced substring
extraction, then \texttt{ValueError} with a 500-character excerpt.
The chain never calls the model again, which preserves the same
post-LLM processing path for every modality and provider, supporting
modality isolation (ER1).

Figure~\ref{fig:impl-json-parse} shows the docstring plus stage~1 only:
the design point is the ordered fallback chain (direct parse, fence
strip, brace-balanced substring, then a bounded \texttt{ValueError}
excerpt). Stages~2--4 are visible in the full function in the appendix.

\begin{figure}[ht]
    \centering
    \implscreenshot{Chapters/figures/fig-6-9-json-parse.png}{%
      Tight crop: \texttt{src/pipeline.py} lines 54--67
      (\texttt{\_parse\_llm\_json} docstring and first parse attempt).}
    \caption{JSON hardening contract and stage~1 (\texttt{src/pipeline.py}).}
    \label{fig:impl-json-parse}
\end{figure}

\texttt{run\_pipeline} snapshots \texttt{config.yaml}, the manifest,
and \texttt{prompt.txt} into the run directory, then iterates
documents with resume-by-default (\texttt{extraction.json} skip)
and per-document \texttt{try}/\texttt{except} that writes
\texttt{error.json} on failure. That is the mechanical realisation of
NFR1 (configuration and prompt snapshots), NFR3 (per-document error
isolation), and ER4 (artefacts sufficient to reconstruct the run)
from Chapter~\ref{chap:spec}.

\apxref{C.7}

% ------------------------------------------------------------
\section{Normalisation}
\label{sec:impl-normalisation}

\texttt{FIELD\_NORMALISERS} maps each schema field name to a
function. \texttt{normalise\_extraction} walks that map so adding a
field touches one Pydantic attribute and one table entry, not a
switch statement scattered across three files.

\texttt{\_parse\_numeric\_string} disambiguates European
(\texttt{1.234,56}) and US (\texttt{1,234.56}) formats by the
relative order of the last comma and dot, not by locale, because LLM
output does not carry locale metadata. That is a concrete
implementation bet: wrong guess on ambiguous \texttt{12,34} is
possible, but bills in the dataset rarely use two-digit
``thousands'' groups without a third group.

Figure~\ref{fig:impl-normalisers} shows the dispatch table only: the
design point is that new fields add one Pydantic attribute and one
table row instead of a scattered \texttt{if} chain. The EU/US numeric
parser and per-field normaliser bodies are described above and listed
in full in the appendix.

\begin{figure}[ht]
    \centering
    \implscreenshot{Chapters/figures/fig-6-10-normalisers.png}{%
      Tight crop: \texttt{src/normalisation.py} lines 317--332
      (\texttt{FIELD\_NORMALISERS}).}
    \caption{Dispatch-table normalisation (\texttt{src/normalisation.py}).}
    \label{fig:impl-normalisers}
\end{figure}

\apxref{C.8}

% ------------------------------------------------------------
\section{Evaluation and diagnosis}
\label{sec:impl-evaldiag}

\texttt{compare\_field} resolves the null 2$\times$2 table before
any value comparison, then applies type-specific rules: $\pm$0.01
tolerance on money and consumption floats, exact string match on
ISO dates after normalisation, exact match on other strings with a
Levenshtein ratio stored alongside for inspection only (pass or
fail remains exact match).

Figure~\ref{fig:impl-compare-field} shows \texttt{compare\_field} from
its signature through line~124: null flags, the full 2$\times$2 null
table, then the float and date arms and the opening of the string
\texttt{else}. The figure uses a compact layout for the null returns;
the string-comparison body, the \texttt{category} assignment, and the
final \texttt{return FieldResult} sit below the crop.

\begin{figure}[ht]
    \centering
    \implscreenshot{Chapters/figures/fig-6-11-compare-field.png}{%
      Tight crop: \texttt{src/evaluation/metrics.py} lines 70--124
      (\texttt{compare\_field}: null table and typed branches).}
    \caption{Null-aware field comparison (\texttt{src/evaluation/metrics.py}).}
    \label{fig:impl-compare-field}
\end{figure}

\apxref{C.9}

\texttt{diagnose\_document} reads \texttt{ocr\_text.txt} for
\texttt{ocr\_text} mode and applies the substring heuristic from
Chapter~\ref{chap:design}, which is the automated half of ER5
(failure attribution in OCR-mediated conditions). Vision mode labels
every non-correct field as \texttt{vision\_model\_failure}. The
implementation writes \texttt{diagnosis.json} only when the CLI passes
\texttt{--diagnose}, so batch runs stay cheap by default while FR6
remains available on demand.

Figure~\ref{fig:impl-diagnosis} shows the OCR-mode loop only: the
condition \texttt{fr.ground\_truth is None or \_value\_in\_text(...)}
means a missing ground-truth string skips substring search and is
labelled \texttt{llm\_extraction\_failure} (hallucination), while a
present ground truth is checked against \texttt{ocr\_text.txt} to
split OCR loss from LLM mis-read. Vision-mode diagnosis and JSON
serialisation are not in the crop.

\begin{figure}[ht]
    \centering
    \implscreenshot{Chapters/figures/fig-6-12-diagnosis.png}{%
      Tight crop: \texttt{src/evaluation/diagnosis.py} lines 62--73
      (OCR-mode \texttt{failure\_type} assignment).}
    \caption{OCR-mode failure attribution (\texttt{src/evaluation/diagnosis.py}).}
    \label{fig:impl-diagnosis}
\end{figure}

\apxref{C.10}

% ------------------------------------------------------------
\section{Deviations from the design}
\label{sec:impl-deviations}

Three deviations were anticipated in
\S\ref{sec:design-maintenance}; all three appear in the code paths
above. First, OpenAI and Anthropic use different APIs and different
output guarantees; that asymmetry is visible in
Figures~\ref{fig:impl-openai-text} and~\ref{fig:impl-anthropic-text}.
Second, wall-clock latency is always measured with \texttt{Timer}
around the pipeline stages; provider-reported latency fields are
not trusted for aggregation. Third, \texttt{TesseractOCR} catches
\texttt{TesseractError} and returns empty text so a missing
language pack degrades to an OCR-side failure mode that diagnosis
can label.

Two additional deviations surfaced during implementation. Fourth,
\texttt{compare\_field} records a Levenshtein ratio for string
fields but still classifies pass or fail on exact equality only.
The design text in Chapter~\ref{chap:design} mentioned similarity
for inspection; the code matches that: the ratio is persisted in
\texttt{evaluation.json} for analysis but does not change the
four-class label. Fifth, per-model USD pricing for
\texttt{estimate\_cost} lives in a hard-coded dict inside
\texttt{src/performance.py} rather than in \texttt{config.yaml}.
The trade is reproducibility from a single checkout: a marker who
checks out the same commit sees the same price table that produced
the run's \texttt{estimated\_cost\_usd}. Moving prices to YAML
would be a small refactor if billing rates need frequent edits.

None of these five items change the factorial structure, the null
taxonomy, or the modality-isolation boundary that the evaluation
chapter will test.

% ------------------------------------------------------------
\section{Testing}
\label{sec:impl-testing}

Tests are organised one file per module, with class names that read
as behaviour contracts (\texttt{TestNormaliseDateShouldParseKnownFormats},
\texttt{TestMissingColumnsShouldRaise}, and similar).
There are \texttt{531} \texttt{def test\_*} functions across the
twelve files; the densest files are
\texttt{test\_evaluation.py}, \texttt{test\_pipeline.py},
\texttt{test\_normalisation.py}, and \texttt{test\_dataset\_loader.py}.

Figure~\ref{fig:impl-tests} is a two-panel composite: the left crop
shows the parametrised date table (positive coverage in one place);
the right crop shows assertions on \texttt{error.json} after bad LLM
JSON (pipeline resilience). Other test modules are referenced only in
prose.

\begin{figure}[ht]
    \centering
    \implscreenshot{Chapters/figures/fig-6-13-tests.png}{%
      Two crops: \texttt{tests/test\_normalisation.py} lines 34--45;
      \texttt{tests/test\_pipeline.py} lines 576--582.}
    \caption{Representative tests: parametrised dates and
    \texttt{error.json} on parse failure.}
    \label{fig:impl-tests}
\end{figure}

\apxref{C.11}

% ------------------------------------------------------------
\section{Running a condition end to end}
\label{sec:impl-running}

The surface area is four Click commands:
\texttt{run}, \texttt{evaluate}, \texttt{compare}, \texttt{report}.
\texttt{run} deep-merges CLI overrides into YAML, executes
\texttt{run\_pipeline}, and prints the run directory path.
\texttt{evaluate} writes \texttt{evaluation.json};
\texttt{--diagnose} adds \texttt{diagnosis.json} per document.
\texttt{compare} builds \texttt{comparison.json} for the four cells
of the $2 \times 2$ design. \texttt{report} reads
\texttt{comparison.json} and writes the six chart PNGs plus
\texttt{summary.txt} described in Chapter~\ref{chap:spec}.

Figure~\ref{fig:impl-run-cli} shows terminal output from
\texttt{python cli.py run ...} or, if API keys are not configured,
\texttt{python cli.py --help}. The design point is a single Click group
with four subcommands; option wiring and imports are not visible in a
terminal crop.

\begin{figure}[ht]
    \centering
    \implscreenshot{Chapters/figures/fig-6-14-run-cli.png}{%
      Screenshot: terminal, \texttt{python cli.py run ...} (or
      \texttt{python cli.py --help}).}
    \caption{CLI execution: run pipeline or command help
    (\texttt{cli.py}).}
    \label{fig:impl-run-cli}
\end{figure}

\apxref{C.12}

% ------------------------------------------------------------
\section{Dependency and maintenance recap}
\label{sec:impl-maintenance}

Runtime Python dependencies are listed in \texttt{requirements.txt}:
\texttt{pydantic}, \texttt{pyyaml}, \texttt{pdf2image},
\texttt{Pillow}, \texttt{pytesseract}, \texttt{openai},
\texttt{anthropic}, \texttt{python-Levenshtein}, \texttt{click},
\texttt{matplotlib}, \texttt{seaborn}, \texttt{pandas},
\texttt{psutil}, \texttt{pytest}, \texttt{python-dotenv}. System
binaries still required on the host are Poppler (for
\texttt{pdf2image}) and Tesseract with language packs
\texttt{eng}, \texttt{deu}, \texttt{fra}, \texttt{ita}. Extensibility
remains the three one-touch additions from
Chapter~\ref{chap:design}: new provider module plus registry line;
new \texttt{OCREngine} subclass; new schema field plus normaliser
entry.

% ------------------------------------------------------------
\section{Summary and forward bridge}
\label{sec:impl-summary}

This chapter mapped the ten design decisions of
Figure~\ref{fig:design-decisions} to concrete modules and showed the
five bounded deviations in \S\ref{sec:impl-deviations}. The
implementation is about \texttt{2,700} lines of runtime Python and
\texttt{531} tests, with disk artefacts and CLI commands
that make a full condition reproducible from a checkout plus a run
directory snapshot.

Two experimental requirements from Chapter~\ref{chap:spec} anchor
the narrative above. ER1 (modality isolation: same reasoning engine,
prompt template, and post-processing across modalities within a
provider row) is realised by the single shared
\texttt{\_process\_document} path after rasterisation
(\S\ref{sec:impl-pipeline}), the shared \texttt{\_parse\_llm\_json}
and normalisation stages, and the identical \texttt{BillExtraction}
contract for both providers (\S\S\ref{sec:impl-schema}--\ref{sec:impl-normalisation}).
ER4 (reproducibility: every run archives configuration, manifest,
prompt, and per-document artefacts) is realised by
\texttt{run\_pipeline}'s directory snapshots and outputs
(\S\ref{sec:impl-pipeline}), which \texttt{tests/test\_pipeline.py}
exercises together with resume and error paths.

Chapter~\ref{chap:evaluation} executes the $2 \times 2$ experiment,
reports accuracy, cost, and latency by condition, slices results by
language and utility type, and returns to the rejected alternatives
in Figure~\ref{fig:design-decisions} to compare predicted outcomes
with observed ones.
.
%  Figures fig-6-1 .. fig-6-14 are user screenshots; each is wrapped
%  in \IfFileExists so the document compiles before PNGs are added.
% ============================================================

\chapter{Implementation}
\label{chap:impl}

% Helper: show image if present, else visible placeholder box.
\newcommand{\implscreenshot}[2]{%
  \IfFileExists{#1}{%
    \includegraphics[width=\linewidth,keepaspectratio]{#1}}{%
    \fbox{\begin{minipage}{0.96\linewidth}\centering\footnotesize
    \texttt{#1}\\[0.35em]#2\end{minipage}}}}

% Pointer to future appendix material (B.x or C.x). Replace macro body
% when appendices are written so all chapter pointers update together.
\newcommand{\apxref}[1]{\textit{(see Appendix~#1 when added.)}}

% ------------------------------------------------------------
\section{From design to implementation}
\label{sec:impl-bridge}

Chapter~\ref{chap:design} fixed what must be built and why each
major choice beat its alternative (Figure~\ref{fig:design-decisions}).
This chapter shows where that design lives in code, what
non-trivial implementation choices were made while wiring it, and
where the running system diverged from the plan
(\S\ref{sec:impl-deviations}). Every row of
Figure~\ref{fig:design-decisions} maps to at least one module or
identifiable block listed in \S\ref{sec:impl-layout}; the four
hardest blocks (provider asymmetry, JSON hardening, null-aware
comparison, failure diagnosis) each get a screenshot in
\S\S\ref{sec:impl-provider}--\ref{sec:impl-evaldiag}.

% ------------------------------------------------------------
\section{Repository layout}
\label{sec:impl-layout}

The runtime is a single Python package under \texttt{src/} plus a
thin \texttt{cli.py} entry point. Together they are about
\texttt{2,700} lines. Evaluation, LLM adapters, and OCR sit in
subpackages so the pipeline imports a stable surface
(\texttt{DatasetLoader}, \texttt{get\_provider}, \texttt{run\_pipeline},
\texttt{evaluate\_run}, \texttt{compare\_runs},
\texttt{generate\_report}). The test tree under \texttt{tests/}
contains \texttt{531} test functions across twelve files
(\S\ref{sec:impl-testing}). Configuration and the zero-shot prompt
live outside the package at \texttt{config/default.yaml} and
\texttt{prompts/extraction\_v1.txt} so a run can be replayed from
artefacts without re-reading the source.

Figure~\ref{fig:impl-tree} evidences a single package root with no
parallel ``evaluation copy'' of the schema: the same
\texttt{BillExtraction} type is imported by the pipeline, both
providers, and the evaluator. Individual file purposes and line
counts are not readable from a tree alone.

\begin{figure}[ht]
    \centering
    \implscreenshot{Chapters/figures/fig-6-1-project-tree.png}{%
      Screenshot: project tree (\texttt{src/}, \texttt{tests/},
      \texttt{config/}, \texttt{prompts/}, \texttt{cli.py}).}
    \caption{Repository layout. Runtime code under \texttt{src/};
    tests mirror module names.}
    \label{fig:impl-tree}
\end{figure}

A per-module inventory (path, purpose, key symbols, approximate
line count) will be tabulated when the appendices are finalised.
\apxref{B.1}

% ------------------------------------------------------------
\section{Data contract}
\label{sec:impl-schema}

\texttt{BillExtraction} is a Pydantic v2 \texttt{BaseModel} with
twelve optional fields and a class-level \texttt{SCHEMA\_FIELDS}
list that fixes iteration order for evaluation and reporting.
\texttt{schema\_description()} builds the bullet list injected into
the prompt at \texttt{\{schema\_description\}}; the same strings
appear nowhere else, so a field rename is a single edit.

The first implementation attempt used Pydantic v1 because the
course template did not pin a version. OpenAI's Structured Outputs
call (\texttt{text\_format=BillExtraction}) failed to emit a valid
JSON schema from the model class. Pinning Pydantic
\texttt{\textgreater=2.0} in \texttt{requirements.txt} and
migrating the model to v2 metadata fixed the failure in one step.
The alternative, dropping Structured Outputs and parsing JSON by
hand for OpenAI too, was rejected because it would remove the
strongest guarantee on the modality comparison's OpenAI arm.

Figure~\ref{fig:impl-schema} shows only the canonical iteration order
\texttt{SCHEMA\_FIELDS}: the design point is that evaluation, reporting,
and the prompt builder all walk the same ordered list. The twelve
\texttt{Optional} field declarations and \texttt{schema\_description()}
live in the same file but are omitted here to keep the crop short.

\begin{figure}[ht]
    \centering
    \implscreenshot{Chapters/figures/fig-6-2-schema.png}{%
      Tight crop: \texttt{src/schema.py} lines 33--46
      (\texttt{SCHEMA\_FIELDS} list).}
    \caption{Canonical field order for iteration (\texttt{src/schema.py}).}
    \label{fig:impl-schema}
\end{figure}

\apxref{C.1}

% ------------------------------------------------------------
\section{Manifest-driven loading}
\label{sec:impl-loader}

\texttt{DatasetLoader.load\_and\_validate} chains six steps: read
CSV, assert required columns, reject duplicate
\texttt{document\_id}, validate cell values (language, utility type,
status, integer \texttt{page\_count}), filter to
\texttt{active} rows with \texttt{annotated} and \texttt{verified}
true, resolve PDF and ground-truth paths, then assert both files
exist. An early prototype collapsed validation into one function
that returned only the first error string. That made large
manifests painful to debug because a row with a bad language code
and a missing PDF looked like a single generic failure. Splitting
the chain keeps each failure mode attributable to one stage, which
is what reproducibility audits need when a row is rejected.

Figure~\ref{fig:impl-loader} shows the ordered call chain inside
\texttt{load\_and\_validate}: the design point is staged validation so
each failure mode maps to one step. Helper implementations
(\texttt{\_read\_csv}, \texttt{\_validate\_columns}, and the rest) are
not shown in the crop.

\begin{figure}[ht]
    \centering
    \implscreenshot{Chapters/figures/fig-6-3-loader.png}{%
      Tight crop: \texttt{src/dataset\_loader.py} lines 55--61
      (sequential validation chain).}
    \caption{Staged manifest validation (\texttt{src/dataset\_loader.py}).}
    \label{fig:impl-loader}
\end{figure}

\apxref{C.2}

% ------------------------------------------------------------
\section{Provider strategy in code}
\label{sec:impl-provider}

\subsection{Abstract contract}

\texttt{LLMProvider} defines \texttt{extract\_from\_text},
\texttt{extract\_from\_image}, and \texttt{get\_model\_id}. Both
extraction methods return the same dict shape so
\texttt{pipeline.py} never branches on provider name. The prompt is
always the rendered instruction; OCR text or image bytes are
passed as the user payload only, which keeps token counts stable
when the instruction is long.

Figure~\ref{fig:impl-llm-abc} shows the text-mode abstract signature and
its return-shape contract: the design point is that both modalities
must return the same dict keys so \texttt{pipeline.py} stays
provider-agnostic. The vision abstract and shared helpers are omitted
from the crop.

\begin{figure}[ht]
    \centering
    \implscreenshot{Chapters/figures/fig-6-4-llm-abc.png}{%
      Tight crop: \texttt{src/llm/base.py} lines 45--58
      (\texttt{extract\_from\_text} signature and Returns block).}
    \caption{Provider contract for text mode (\texttt{src/llm/base.py}).}
    \label{fig:impl-llm-abc}
\end{figure}

\apxref{C.3}

\subsection{OpenAI: Responses API and Structured Outputs}

\texttt{OpenAIProvider} calls \texttt{responses.parse} with
\texttt{text\_format=BillExtraction}. The API returns text that
already matches the schema; the pipeline still runs
\texttt{\_parse\_llm\_json} for uniformity, but the happy path is a
single \texttt{json.loads} on a clean string. This is the direct
code realisation of the ``Structured output handling'' row in
Figure~\ref{fig:design-decisions}.

Figure~\ref{fig:impl-openai-text} isolates the Structured Outputs call:
\texttt{instructions=prompt} and a user message carrying
\texttt{ocr\_text} split rules from data, matching
Chapter~\ref{chap:design}. Latency wrapping and result assembly are
not shown.

\begin{figure}[ht]
    \centering
    \implscreenshot{Chapters/figures/fig-6-5-openai-text.png}{%
      Tight crop: \texttt{src/llm/openai\_provider.py} lines 63--74
      (\texttt{responses.parse} with \texttt{text\_format=BillExtraction}).}
    \caption{OpenAI text mode with Structured Outputs
    (\texttt{src/llm/openai\_provider.py}).}
    \label{fig:impl-openai-text}
\end{figure}

\apxref{C.4}

\subsection{Anthropic: Messages API and fence stripping}

\texttt{AnthropicProvider} uses \texttt{messages.create} with the
same rendered prompt as the system message and OCR text or image
blocks as the user message. There is no Structured Outputs hook
for this schema, so \texttt{\_strip\_json\_fencing} removes
\texttt{```json ... ```} wrappers before the shared parser runs;
that helper is small enough to describe in prose only. Provider-specific
cleaning stops here; anything that survives a malformed fence still
hits the four-stage chain in \S\ref{sec:impl-pipeline}.

Figure~\ref{fig:impl-anthropic-text} shows the text-mode API call only:
\texttt{system=prompt} mirrors OpenAI's split of instruction versus
payload. Vision mode and \texttt{\_build\_result} are omitted from the
crop.

\begin{figure}[ht]
    \centering
    \implscreenshot{Chapters/figures/fig-6-6-anthropic-text.png}{%
      Tight crop: \texttt{src/llm/anthropic\_provider.py} lines 75--88
      (\texttt{messages.create} for text mode).}
    \caption{Anthropic text mode (\texttt{src/llm/anthropic\_provider.py}).}
    \label{fig:impl-anthropic-text}
\end{figure}

\apxref{C.5}

\subsection{Registry and lazy import}

\texttt{registry.py} maps string keys to dotted class paths and
imports the class only when \texttt{get\_provider} runs. That means
a host configured for OpenAI never imports \texttt{anthropic}, which
matters when only one SDK is installed. Adding a third provider is
one new module plus one line in \texttt{\_PROVIDERS}.

Figure~\ref{fig:impl-registry} shows only the string-to-class map: the
design point is that adding a provider is one dictionary entry plus
one module. \texttt{get\_provider}, lazy import, and constructor wiring
are left to the appendix listing.

\begin{figure}[ht]
    \centering
    \implscreenshot{Chapters/figures/fig-6-7-registry.png}{%
      Tight crop: \texttt{src/llm/registry.py} lines 16--19
      (\texttt{\_PROVIDERS}).}
    \caption{Provider registry map (\texttt{src/llm/registry.py}).}
    \label{fig:impl-registry}
\end{figure}

\apxref{C.6}

% ------------------------------------------------------------
\section{Pipeline orchestration}
\label{sec:impl-pipeline}

\texttt{\_process\_document} wraps each stage in
\texttt{Timer} from \texttt{src/performance.py}. The keys written to
\texttt{timings.json} (\texttt{pdf\_to\_images\_ms},
\texttt{ocr\_ms}, \texttt{llm\_call\_ms},
\texttt{parse\_normalise\_ms}, \texttt{total\_ms}) match the design
sequence diagram (Figure~\ref{fig:design-seq}). Vision mode uses a
zero-millisecond placeholder timer for OCR so downstream reporting
can still sum columns without branching on key presence.

Figure~\ref{fig:impl-pipeline-core} shows both arms of the modality
split: OCR-then-\texttt{extract\_from\_text} versus vision-then
\texttt{extract\_from\_image}, including the zero-millisecond OCR timer
placeholder in vision mode. Stage~4 parse, timings assembly, artefact
writes, and cost estimation sit below the crop.

\begin{figure}[ht]
    \centering
    \implscreenshot{Chapters/figures/fig-6-8-pipeline-core.png}{%
      Tight crop: \texttt{src/pipeline.py} lines 133--150
      (\texttt{if mode == "ocr\_text"} and \texttt{else} vision branch).}
    \caption{OCR and vision modality branches inside
    \texttt{\_process\_document} (\texttt{src/pipeline.py}).}
    \label{fig:impl-pipeline-core}
\end{figure}

\texttt{\_parse\_llm\_json} implements the four-stage hardening
chain from Figure~\ref{fig:design-parse}: direct
\texttt{json.loads}, fenced-block strip, brace-balanced substring
extraction, then \texttt{ValueError} with a 500-character excerpt.
The chain never calls the model again, which preserves the same
post-LLM processing path for every modality and provider, supporting
modality isolation (ER1).

Figure~\ref{fig:impl-json-parse} shows the docstring plus stage~1 only:
the design point is the ordered fallback chain (direct parse, fence
strip, brace-balanced substring, then a bounded \texttt{ValueError}
excerpt). Stages~2--4 are visible in the full function in the appendix.

\begin{figure}[ht]
    \centering
    \implscreenshot{Chapters/figures/fig-6-9-json-parse.png}{%
      Tight crop: \texttt{src/pipeline.py} lines 54--67
      (\texttt{\_parse\_llm\_json} docstring and first parse attempt).}
    \caption{JSON hardening contract and stage~1 (\texttt{src/pipeline.py}).}
    \label{fig:impl-json-parse}
\end{figure}

\texttt{run\_pipeline} snapshots \texttt{config.yaml}, the manifest,
and \texttt{prompt.txt} into the run directory, then iterates
documents with resume-by-default (\texttt{extraction.json} skip)
and per-document \texttt{try}/\texttt{except} that writes
\texttt{error.json} on failure. That is the mechanical realisation of
NFR1 (configuration and prompt snapshots), NFR3 (per-document error
isolation), and ER4 (artefacts sufficient to reconstruct the run)
from Chapter~\ref{chap:spec}.

\apxref{C.7}

% ------------------------------------------------------------
\section{Normalisation}
\label{sec:impl-normalisation}

\texttt{FIELD\_NORMALISERS} maps each schema field name to a
function. \texttt{normalise\_extraction} walks that map so adding a
field touches one Pydantic attribute and one table entry, not a
switch statement scattered across three files.

\texttt{\_parse\_numeric\_string} disambiguates European
(\texttt{1.234,56}) and US (\texttt{1,234.56}) formats by the
relative order of the last comma and dot, not by locale, because LLM
output does not carry locale metadata. That is a concrete
implementation bet: wrong guess on ambiguous \texttt{12,34} is
possible, but bills in the dataset rarely use two-digit
``thousands'' groups without a third group.

Figure~\ref{fig:impl-normalisers} shows the dispatch table only: the
design point is that new fields add one Pydantic attribute and one
table row instead of a scattered \texttt{if} chain. The EU/US numeric
parser and per-field normaliser bodies are described above and listed
in full in the appendix.

\begin{figure}[ht]
    \centering
    \implscreenshot{Chapters/figures/fig-6-10-normalisers.png}{%
      Tight crop: \texttt{src/normalisation.py} lines 317--332
      (\texttt{FIELD\_NORMALISERS}).}
    \caption{Dispatch-table normalisation (\texttt{src/normalisation.py}).}
    \label{fig:impl-normalisers}
\end{figure}

\apxref{C.8}

% ------------------------------------------------------------
\section{Evaluation and diagnosis}
\label{sec:impl-evaldiag}

\texttt{compare\_field} resolves the null 2$\times$2 table before
any value comparison, then applies type-specific rules: $\pm$0.01
tolerance on money and consumption floats, exact string match on
ISO dates after normalisation, exact match on other strings with a
Levenshtein ratio stored alongside for inspection only (pass or
fail remains exact match).

Figure~\ref{fig:impl-compare-field} shows \texttt{compare\_field} from
its signature through line~124: null flags, the full 2$\times$2 null
table, then the float and date arms and the opening of the string
\texttt{else}. The figure uses a compact layout for the null returns;
the string-comparison body, the \texttt{category} assignment, and the
final \texttt{return FieldResult} sit below the crop.

\begin{figure}[ht]
    \centering
    \implscreenshot{Chapters/figures/fig-6-11-compare-field.png}{%
      Tight crop: \texttt{src/evaluation/metrics.py} lines 70--124
      (\texttt{compare\_field}: null table and typed branches).}
    \caption{Null-aware field comparison (\texttt{src/evaluation/metrics.py}).}
    \label{fig:impl-compare-field}
\end{figure}

\apxref{C.9}

\texttt{diagnose\_document} reads \texttt{ocr\_text.txt} for
\texttt{ocr\_text} mode and applies the substring heuristic from
Chapter~\ref{chap:design}, which is the automated half of ER5
(failure attribution in OCR-mediated conditions). Vision mode labels
every non-correct field as \texttt{vision\_model\_failure}. The
implementation writes \texttt{diagnosis.json} only when the CLI passes
\texttt{--diagnose}, so batch runs stay cheap by default while FR6
remains available on demand.

Figure~\ref{fig:impl-diagnosis} shows the OCR-mode loop only: the
condition \texttt{fr.ground\_truth is None or \_value\_in\_text(...)}
means a missing ground-truth string skips substring search and is
labelled \texttt{llm\_extraction\_failure} (hallucination), while a
present ground truth is checked against \texttt{ocr\_text.txt} to
split OCR loss from LLM mis-read. Vision-mode diagnosis and JSON
serialisation are not in the crop.

\begin{figure}[ht]
    \centering
    \implscreenshot{Chapters/figures/fig-6-12-diagnosis.png}{%
      Tight crop: \texttt{src/evaluation/diagnosis.py} lines 62--73
      (OCR-mode \texttt{failure\_type} assignment).}
    \caption{OCR-mode failure attribution (\texttt{src/evaluation/diagnosis.py}).}
    \label{fig:impl-diagnosis}
\end{figure}

\apxref{C.10}

% ------------------------------------------------------------
\section{Deviations from the design}
\label{sec:impl-deviations}

Three deviations were anticipated in
\S\ref{sec:design-maintenance}; all three appear in the code paths
above. First, OpenAI and Anthropic use different APIs and different
output guarantees; that asymmetry is visible in
Figures~\ref{fig:impl-openai-text} and~\ref{fig:impl-anthropic-text}.
Second, wall-clock latency is always measured with \texttt{Timer}
around the pipeline stages; provider-reported latency fields are
not trusted for aggregation. Third, \texttt{TesseractOCR} catches
\texttt{TesseractError} and returns empty text so a missing
language pack degrades to an OCR-side failure mode that diagnosis
can label.

Two additional deviations surfaced during implementation. Fourth,
\texttt{compare\_field} records a Levenshtein ratio for string
fields but still classifies pass or fail on exact equality only.
The design text in Chapter~\ref{chap:design} mentioned similarity
for inspection; the code matches that: the ratio is persisted in
\texttt{evaluation.json} for analysis but does not change the
four-class label. Fifth, per-model USD pricing for
\texttt{estimate\_cost} lives in a hard-coded dict inside
\texttt{src/performance.py} rather than in \texttt{config.yaml}.
The trade is reproducibility from a single checkout: a marker who
checks out the same commit sees the same price table that produced
the run's \texttt{estimated\_cost\_usd}. Moving prices to YAML
would be a small refactor if billing rates need frequent edits.

None of these five items change the factorial structure, the null
taxonomy, or the modality-isolation boundary that the evaluation
chapter will test.

% ------------------------------------------------------------
\section{Testing}
\label{sec:impl-testing}

Tests are organised one file per module, with class names that read
as behaviour contracts (\texttt{TestNormaliseDateShouldParseKnownFormats},
\texttt{TestMissingColumnsShouldRaise}, and similar).
There are \texttt{531} \texttt{def test\_*} functions across the
twelve files; the densest files are
\texttt{test\_evaluation.py}, \texttt{test\_pipeline.py},
\texttt{test\_normalisation.py}, and \texttt{test\_dataset\_loader.py}.

Figure~\ref{fig:impl-tests} is a two-panel composite: the left crop
shows the parametrised date table (positive coverage in one place);
the right crop shows assertions on \texttt{error.json} after bad LLM
JSON (pipeline resilience). Other test modules are referenced only in
prose.

\begin{figure}[ht]
    \centering
    \implscreenshot{Chapters/figures/fig-6-13-tests.png}{%
      Two crops: \texttt{tests/test\_normalisation.py} lines 34--45;
      \texttt{tests/test\_pipeline.py} lines 576--582.}
    \caption{Representative tests: parametrised dates and
    \texttt{error.json} on parse failure.}
    \label{fig:impl-tests}
\end{figure}

\apxref{C.11}

% ------------------------------------------------------------
\section{Running a condition end to end}
\label{sec:impl-running}

The surface area is four Click commands:
\texttt{run}, \texttt{evaluate}, \texttt{compare}, \texttt{report}.
\texttt{run} deep-merges CLI overrides into YAML, executes
\texttt{run\_pipeline}, and prints the run directory path.
\texttt{evaluate} writes \texttt{evaluation.json};
\texttt{--diagnose} adds \texttt{diagnosis.json} per document.
\texttt{compare} builds \texttt{comparison.json} for the four cells
of the $2 \times 2$ design. \texttt{report} reads
\texttt{comparison.json} and writes the six chart PNGs plus
\texttt{summary.txt} described in Chapter~\ref{chap:spec}.

Figure~\ref{fig:impl-run-cli} shows terminal output from
\texttt{python cli.py run ...} or, if API keys are not configured,
\texttt{python cli.py --help}. The design point is a single Click group
with four subcommands; option wiring and imports are not visible in a
terminal crop.

\begin{figure}[ht]
    \centering
    \implscreenshot{Chapters/figures/fig-6-14-run-cli.png}{%
      Screenshot: terminal, \texttt{python cli.py run ...} (or
      \texttt{python cli.py --help}).}
    \caption{CLI execution: run pipeline or command help
    (\texttt{cli.py}).}
    \label{fig:impl-run-cli}
\end{figure}

\apxref{C.12}

% ------------------------------------------------------------
\section{Dependency and maintenance recap}
\label{sec:impl-maintenance}

Runtime Python dependencies are listed in \texttt{requirements.txt}:
\texttt{pydantic}, \texttt{pyyaml}, \texttt{pdf2image},
\texttt{Pillow}, \texttt{pytesseract}, \texttt{openai},
\texttt{anthropic}, \texttt{python-Levenshtein}, \texttt{click},
\texttt{matplotlib}, \texttt{seaborn}, \texttt{pandas},
\texttt{psutil}, \texttt{pytest}, \texttt{python-dotenv}. System
binaries still required on the host are Poppler (for
\texttt{pdf2image}) and Tesseract with language packs
\texttt{eng}, \texttt{deu}, \texttt{fra}, \texttt{ita}. Extensibility
remains the three one-touch additions from
Chapter~\ref{chap:design}: new provider module plus registry line;
new \texttt{OCREngine} subclass; new schema field plus normaliser
entry.

% ------------------------------------------------------------
\section{Summary and forward bridge}
\label{sec:impl-summary}

This chapter mapped the ten design decisions of
Figure~\ref{fig:design-decisions} to concrete modules and showed the
five bounded deviations in \S\ref{sec:impl-deviations}. The
implementation is about \texttt{2,700} lines of runtime Python and
\texttt{531} tests, with disk artefacts and CLI commands
that make a full condition reproducible from a checkout plus a run
directory snapshot.

Two experimental requirements from Chapter~\ref{chap:spec} anchor
the narrative above. ER1 (modality isolation: same reasoning engine,
prompt template, and post-processing across modalities within a
provider row) is realised by the single shared
\texttt{\_process\_document} path after rasterisation
(\S\ref{sec:impl-pipeline}), the shared \texttt{\_parse\_llm\_json}
and normalisation stages, and the identical \texttt{BillExtraction}
contract for both providers (\S\S\ref{sec:impl-schema}--\ref{sec:impl-normalisation}).
ER4 (reproducibility: every run archives configuration, manifest,
prompt, and per-document artefacts) is realised by
\texttt{run\_pipeline}'s directory snapshots and outputs
(\S\ref{sec:impl-pipeline}), which \texttt{tests/test\_pipeline.py}
exercises together with resume and error paths.

Chapter~\ref{chap:evaluation} executes the $2 \times 2$ experiment,
reports accuracy, cost, and latency by condition, slices results by
language and utility type, and returns to the rejected alternatives
in Figure~\ref{fig:design-decisions} to compare predicted outcomes
with observed ones.
